% =========================================================
% Proof Stub: Direct Comparison Test
% Source: volume-iii/analysis/sequences/notes/series/notes-series-tests.tex
% Mode: standard
% =========================================================

\clearpage
\phantomsection
\hypertarget{proof-RA-ABB-T-direct-comparison}{}

\subsubsection[Direct Comparison Test]{Proof --- Direct Comparison Test}

\bigskip

\noindent
\textbf{Source.}
Abbott, \emph{Understanding Analysis}; see \texttt{notes-series-tests.tex}.

\vspace{0.75em}

\noindent
\textbf{Goal.}
If $0 \le a_n \le b_n$ eventually and $\sum b_n < \infty$, then $\sum a_n < \infty$.

\vspace{0.75em}

\noindent
\textbf{Logical form.}
\[
  a_n \le b_n \ge 0 \text{ eventually},\; \sum b_n < \infty \Rightarrow \sum a_n < \infty.
\]

\vspace{0.75em}

\noindent
\textbf{Proof strategy.}
Partial sums of $(a_n)$ are bounded above by partial sums of $(b_n)$, which are bounded. The partial sums of $(a_n)$ are also increasing (since $a_n \ge 0$). MCT applies.

\vspace{0.75em}

\noindent
\textbf{Proof.}
\begin{proof}
\Claim If $0 \le a_n \le b_n$ eventually and $\sum b_n < \infty$, then $\sum a_n < \infty$.

\Given 

\Goal 

\Strategy Partial sums of $(a_n)$ are bounded above by partial sums of $(b_n)$, which are bounded. The partial sums of $(a_n)$ are also increasing (since $a_n \ge 0$). MCT applies.

\medskip

% --- let s_n = ∑_{k=1}^n a_k, t_n = ∑_{k=1}^n b_k ---
% --- s_n ≤ t_n (by comparison, up to the eventual start) ---
% --- t_n ≤ ∑b_n < ∞ for all n ---
% --- s_n increasing (a_n ≥ 0) and bounded above ---
% --- MCT: s_n converges ---

\end{proof}

\begin{remark*}[Proof shape]
Bounded increasing partial sums converge by MCT.
\end{remark*}

\begin{remark*}[Dependencies]
\begin{itemize}
  \item Monotone Convergence Theorem.
  \item $a_n \ge 0$ ensures partial sums are increasing.
  \item Comparison $a_n \le b_n$ gives the upper bound.
\end{itemize}
\end{remark*}
